\section{Results \& Analysis}\label{sec:results}

Table~\ref{tab:key_param_comparison} reports posterior means and 95\% quantile-based CI for key parameters across all fitted models. Entries marked \textemdash{} indicate that the parameter is not present in that model specification.

\begin{table}[h!]
\centering
\caption{Posterior mean and 95\% CI for selected parameters across models.}
\label{tab:key_param_comparison}
\begin{tabular}{lccccc}
\hline
Model & $\alpha$ & $\beta_{\mathrm{fn}}$ & $\beta_{\mathrm{fire\_size}}$ & $\alpha_{\text{prov=BC}}$ & $\beta_{\mathrm{dist\_to\_fn}}$ \\
\hline
Minimal & $-3.4_{[-3.5,\,-3.3]}$ & $1.4_{[1.0,\,1.8]}$ & $0.4_{[0.3,\,0.5]}$ & \textemdash{} & \textemdash{} \\
Province & $-4.0_{[-4.7,\,-3.4]}$ & $1.2_{[0.8,\,1.6]}$ & $0.5_{[0.4,\,0.6]}$ & $1.8_{[1.2,\,2.5]}$ & \textemdash{} \\
Full & $-3.9_{[-4.9,\,-2.9]}$ & $-1.0_{[-1.4,\,-0.5]}$ & $0.5_{[0.4,\,0.6]}$ & $1.5_{[0.9,\,2.2]}$ & $-1.1_{[-1.2,\,-1.0]}$ \\
Response & $-4.0_{[-5.3,\,-2.6]}$ & $-1.0_{[-1.4,\,-0.5]}$ & $0.5_{[0.4,\,0.6]}$ & $1.5_{[0.9,\,2.1]}$ & $-1.1_{[-1.2,\,-0.9]}$ \\
Spatial (VI) & $-2.6_{[-2.7,\,-2.5]}$ & $-0.9_{[-1.3,\,-0.5]}$ & $0.5_{[0.5,\,0.6]}$ & $0.4_{[0.2,\,0.6]}$ & $-1.0_{[-1.1,\,-0.9]}$ \\
\hline
\end{tabular}
\end{table}

\textbf{Intercept:} All model broadly agree on the value for $\alpha$, with Minimal a small outlier and Spatial a larger outlier. Notably the 95\% CI for the Spatial model is much tighter than the rest. As a quick validation, I ran the Full model with VI and found a posterior for $\alpha$ of -2.28 with CI $[-2.4, -2.2]$. However, as with the Spatial model, other parameter values were more consistent with the MCMC inference. It may be that the VI approximation for $\alpha$ is a particularly bad estimate of the posterior. From this, I conclude that we can give some weight to the Spatial model, especially if they agree with the other models, but not too much as the VI approximation certainly loses information.\\
\textbf{FN-Indicator \& Distance to FN Reserve:} The most notable result here is the stark change from a positive to negative $\beta_{\text{fn}}$ between the Minimal/Province models and the remaining models. This is likely due to the presence of the $\beta_{\text{dist\_to\_fn}}$ in the final 3 models. However, its value is negative as is the case for $\beta_{\text{fn}}$ which is unsurprising. It is reasonable that the farther the fire is from a reserve, the less likely an evacuation is to be ordered. However, $\beta_{\text{fn}}$ remains negative, suggesting that a fire being on the reserve is roughly associated with a $\exp(-1)=0.37$ reduction in odds, or 67\% less of an evacuation. Note that these results are purely associative and causal.\\
\textbf{Provinces:} The province with the highest $\alpha$ was always British Columbia (BC) with mean $\alpha_{\text{BC}} \approx 1.5$ for the full model, while Quebec (QC) had $\alpha_{\text{QC}} \approx -1.4$. In general, these coefficients suspiciously correlate with the provinces which have the most fires. It may be that provinces with the most wildfires are more experienced in handling evacuations, are more risk-averse as there is greater potential for the wildfire to grow, or that the model learns that some provinces almost never have enough fires large enough to issue an evacuation. Regardless of the cause, this suggests there are confounders missing about the evacuation practices of each province.\\
\textbf{Responses:} This model was intended to test the impact of including response severity in the model on $\beta_{\text{fn}}$. In a sense, this model is unfair since perfect knowledge of the response could include the evacuation decision. However, the model obtains very similar ELPD and parameter means as the Full model. Indeed, $\beta_{\text{fn}}\approx -1$ for both models. However, the only response $\alpha$ with a positive coefficient is the `other', corresponding to Nans in the original data. Meanwhile a `full' response has a mean of $-0.05$ with a standard deviation (SD) of $0.52$ indicating relatively high uncertainty. \\
\textbf{Protection Zones:} The protection zones show a clearer effect. In the Full model, the `nordique' zone has the largest mean of $1.42$ and SD of $0.61$, while the smallest is `monitor' with a mean of $-1.52$.\\
\textbf{Regions \& Spatial Model:} The regional effects vary from $-1.43$ to $1.1$ where the mean of $\sigma_{\text{region}}$ is $3.65$. This model should be able to isolate further confounders like `remoteness', while retaining regularisation using a prior. However, the model obtains a very similar ELPD as the Full model. This may be in part due to the VI method. The value for $\alpha_{\text{BC}}$ decreases meaningfully to a mean of $0.4$. I assume that much of the `BC effect' has instead been captured by the smaller regional factors in the model. Crucially, the value of $\beta_{\text{fn}}$ remains negative in this model.

